\documentclass[10pt]{article}
\usepackage[T1]{fontenc}
\usepackage{lmodern}
\usepackage[margin=22mm]{geometry}
\usepackage{xcolor}
\usepackage{microtype}
\usepackage{enumitem}
\usepackage{hyperref}
\usepackage{xurl}
\usepackage{forest}
usetikzlibrary{arrows.meta}
\usepackage{pdflscape}
\usepackage[most]{tcolorbox}
\usepackage{titlesec}
\usepackage{fancyhdr}
\definecolor{Navy}{HTML}{17324D}
\definecolor{Teal}{HTML}{176B6B}
\definecolor{PaleBlue}{HTML}{EAF2F8}
\definecolor{FileFill}{HTML}{F7F9FA}
\definecolor{DescText}{HTML}{334E5C}
\definecolor{TreeLine}{HTML}{607D8B}
\definecolor{Warn}{HTML}{8A4B08}
\hypersetup{colorlinks=true,linkcolor=Teal,urlcolor=Teal,pdfborder={0 0 0}}
\titleformat{\section}{\Large\bfseries\color{Navy}}{}{0pt}{}[\color{Teal}\titlerule]
\titleformat{\subsection}{\large\bfseries\color{Teal}}{}{0pt}{}
\setlist[itemize]{leftmargin=*,itemsep=.45em,topsep=.45em}
\setlist[enumerate]{leftmargin=*,itemsep=.45em,topsep=.45em}
\setlength{\parindent}{0pt}
\setlength{\parskip}{.65em}
\newcommand{\code}[1]{\texttt{#1}}
\forestset{
 rootnode/.style={draw=Navy,fill=Navy,text=white,rounded corners=2pt,inner sep=5pt},
 filenode/.style={draw=TreeLine,fill=FileFill,rounded corners=2pt,inner sep=3.2pt},
}
\pagestyle{fancy}\fancyhf{}\fancyhead[L]{\small\color{DescText}Repository Structure Inventory}\fancyhead[R]{\small\color{DescText}\leftmark}\fancyfoot[C]{\thepage}
\setlength{\headheight}{14pt}
\newtcolorbox{verificationbox}{breakable,colback=orange!5,colframe=Warn,title={Verification status},fonttitle=\bfseries}
\begin{document}
\begin{titlepage}
\pagecolor{PaleBlue}
\color{Navy}
\vspace*{0.17\textheight}
{\Huge\bfseries Grade 10 — Term 3 Structure\par}
\vspace{1em}
{\Large Repository resource inventory\par}
\vfill
{\large Authoritative conversion of \code{STRUCTURE.md}\par}
\end{titlepage}
\nopagecolor
\tableofcontents
\clearpage

\section{Overview}

This directory contains the source materials for Grade 10 Term 3 probability instruction. The sequence moves from practice on discrete and continuous distributions (criteria C0–C7), through student assessments and their teacher solution sets, to catch-up work and a C8–C10 web/data-modeling project. All printable resources are LaTeX source files; the only local data assets are three CSV datasets supplied for the project. There are no HTML, CSS, JavaScript, compiled PDFs, or build/configuration files in this directory.

The inventory below documents \textbf{32 existing files in 4 directories}. \code{STRUCTURE.md} itself is the inventory document and is intentionally omitted from the tree.

\section{Directory tree}

\begin{tcolorbox}[colback=PaleBlue,colframe=Teal,title={Reading the diagrams}]
The directory tree is divided into one landscape diagram per top-level directory so that every documented entry and description remains legible. Directory nodes use a dark fill; file nodes use a light fill, with descriptions shown in a secondary text color. The inventory root is \code{.}.
\end{tcolorbox}
\begin{landscape}
\subsection{\code{1-practice\_activities/}}
\begin{center}
\begin{forest}
for tree={grow=east, parent anchor=east, child anchor=west, anchor=west, edge={-Latex,draw=TreeLine}, l sep=9mm, s sep=1.2mm, font=\sffamily\fontsize{6.2}{7.2}\selectfont},
[{\textbf{1-practice\_activities/}\\\textcolor{DescText}{Student practice, consolidated worked material, contexts, and a teacher-oriented summary for C0–C10.}}, rootnode, text width=6.0cm
[{\code{G10\_T3\_C0\_practice\_activity.tex}\\\textcolor{DescText}{Student-facing worked practice on discrete PMFs: validity, unknown probabilities, event sums, graphs, and contextual interpretations.}}, filenode, text width=18.2cm]
[{\code{G10\_T3\_C1C2C3\_practice\_activity\_context.tex}\\\textcolor{DescText}{Student-facing contextual set joining C1 density validity, C2 uniform mixtures, and C3 triangular investment models.}}, filenode, text width=18.2cm]
[{\code{G10\_T3\_C1\_practice\_activity.tex}\\\textcolor{DescText}{Student-facing C1 progression from geometric area review to validating and using uniform and linear PDFs, including parameterized domains.}}, filenode, text width=18.2cm]
[{\code{G10\_T3\_C2\_practice\_activity.tex}\\\textcolor{DescText}{Student-facing C2 practice on uniform and piecewise-uniform PDFs, unknown constants/endpoints, complements, and multi-interval models.}}, filenode, text width=18.2cm]
[{\code{G10\_T3\_C3\_practice\_activity.tex}\\\textcolor{DescText}{Student-facing C3 practice using geometry—not integration—for linear, triangular, piecewise-linear, and parameterized densities.}}, filenode, text width=18.2cm]
[{\code{G10\_T3\_C4\_practice\_activity.tex}\\\textcolor{DescText}{Student-facing C4 worked comparisons of mean, variance, and standard deviation for uniform, triangular, and normal distributions.}}, filenode, text width=18.2cm]
[{\code{G10\_T3\_C5\_practice\_activity.tex}\\\textcolor{DescText}{Student-facing C5 guide and nine worked problem types for reading the standard-normal table and combining tails/intervals.}}, filenode, text width=18.2cm]
[{\code{G10\_T3\_C6\_practice\_activity.tex}\\\textcolor{DescText}{Student-facing C6 contextual normal problems organized by left-tail, right-tail, central-band, and two-region cases after standardization.}}, filenode, text width=18.2cm]
[{\code{G10\_T3\_C7\_practice\_activity.tex}\\\textcolor{DescText}{Student-facing C7 inverse-normal applications for upper/lower percentiles and symmetric central intervals in real contexts.}}, filenode, text width=18.2cm]
[{\code{G10\_T3\_Contexts.tex}\\\textcolor{DescText}{Standalone explanatory article comparing real-life continuous-PDF models, geometric calculations, assumptions, limitations, and references.}}, filenode, text width=18.2cm]
[{\code{G10\_T3\_MATERIAL.tex}\\\textcolor{DescText}{Consolidated worked C1–C10 problem collection; the original edition uses integration in early density solutions.}}, filenode, text width=18.2cm]
[{\code{G10\_T3\_MATERIAL\_V2.tex}\\\textcolor{DescText}{Revised consolidated C1–C10 collection replacing early integration-based work with geometric area reasoning and revising some models.}}, filenode, text width=18.2cm]
[{\code{G10\_T3\_summary.tex}\\\textcolor{DescText}{Teacher-oriented pedagogical summary of the C0–C7 practice files, including progression, strengths, gaps, and suggested improvements.}}, filenode, text width=18.2cm]
]
\end{forest}
\end{center}
\end{landscape}
\begin{landscape}
\subsection{\code{2-learning\_evidences/}}
\begin{center}
\begin{forest}
for tree={grow=east, parent anchor=east, child anchor=west, anchor=west, edge={-Latex,draw=TreeLine}, l sep=9mm, s sep=1.2mm, font=\sffamily\fontsize{6.2}{7.2}\selectfont},
[{\textbf{2-learning\_evidences/}\\\textcolor{DescText}{Lesson assessments, alternate versions, teacher solutions, a computational project brief, and its input datasets.}}, rootnode, text width=6.0cm
[{\code{G10\_T3\_C1C7\_midterm.tex}\\\textcolor{DescText}{Student midterm assessing C1–C7 with original numerical parameters and contextual density/normal-distribution tasks.}}, filenode, text width=18.2cm]
[{\code{G10\_T3\_C1C7\_midtermV2.tex}\\\textcolor{DescText}{Alternate student midterm with revised header/rubric, shifted numerical parameters, and reformatted C5 questions.}}, filenode, text width=18.2cm]
[{\code{G10\_T3\_C1C7\_midterm\_solution.tex}\\\textcolor{DescText}{Teacher solution set for the original C1–C7 midterm, with full derivations, graphs, and a minimal C5–C7 answer section.}}, filenode, text width=18.2cm]
[{\code{G10\_T3\_C1C7\_midterm\_solutionV2.tex}\\\textcolor{DescText}{Teacher solution set matching the V2 midterm's changed endpoints, scenarios, distribution parameters, and answers.}}, filenode, text width=18.2cm]
[{\code{G10\_T3\_L1\_C1C2C3\_exam.tex}\\\textcolor{DescText}{Student Learning Evidence 1 exam: 17 C1–C3 density-validity, uniform, piecewise-uniform, and linear/triangular tasks.}}, filenode, text width=18.2cm]
[{\code{G10\_T3\_L1\_C1C2C3\_exam\_solution.tex}\\\textcolor{DescText}{Teacher worked solutions to the C1–C3 exam, including geometric derivations and plotted density regions.}}, filenode, text width=18.2cm]
[{\code{G10\_T3\_L1\_C4\_exam.tex}\\\textcolor{DescText}{Student C4 Learning Evidence 2 on matching/drawing uniform, triangular, and normal distributions by mean and spread.}}, filenode, text width=18.2cm]
[{\code{G10\_T3\_L1\_C4\_examV2.tex}\\\textcolor{DescText}{Non-credit practice version of the C4 assessment with changed distributions, graph ranges, and answer options.}}, filenode, text width=18.2cm]
[{\code{G10\_T3\_L1\_C4\_examV2\_solution.tex}\\\textcolor{DescText}{Teacher worked answers for the altered C4 V2 practice set, organized problem by problem.}}, filenode, text width=18.2cm]
[{\code{G10\_T3\_L1\_C4\_exam\_solution.tex}\\\textcolor{DescText}{Concise teacher answer key for the original C4 exam, grouped into matching and drawing problems and retaining the exam-style header.}}, filenode, text width=18.2cm]
[{\code{G10\_T3\_L1\_C4\_exam\_solution\_long.tex}\\\textcolor{DescText}{Expanded teacher solution for the original C4 exam with an evaluated-criterion overview and individual explanations for all 18 tasks.}}, filenode, text width=18.2cm]
[{\code{G10\_T3\_L1\_C5\_exam\_solution.tex}\\\textcolor{DescText}{Teacher C5 worked solution set covering six standard-normal table patterns, including symmetry, complements, intervals, and tail unions.}}, filenode, text width=18.2cm]
[{\code{G10\_T3\_L4\_C8C9C10\_activity.tex}\\\textcolor{DescText}{Student Learning Evidence 4 brief requiring an HTML/CSS/JavaScript modeling webpage, AI-development evidence, deployment, and a statistical report.}}, filenode, text width=18.2cm]
[{\code{G10\_T3\_L4\_C8C9C10\_dataset\_1.csv}\\\textcolor{DescText}{C9 test data: 30 measurements labeled as a uniform-like process (12.4–26.3).}}, filenode, text width=18.2cm]
[{\code{G10\_T3\_L4\_C8C9C10\_dataset\_2.csv}\\\textcolor{DescText}{C9 test data: 30 measurements labeled as a triangular-like process (5.2–20.3), concentrated near the middle.}}, filenode, text width=18.2cm]
[{\code{G10\_T3\_L4\_C8C9C10\_dataset\_3.csv}\\\textcolor{DescText}{C9 test data: 30 measurements labeled as a normal-like process (42.1–65.2), centered in the mid-50s.}}, filenode, text width=18.2cm]
]
\end{forest}
\end{center}
\end{landscape}
\begin{landscape}
\subsection{\code{3-catch\_ups/}}
\begin{center}
\begin{forest}
for tree={grow=east, parent anchor=east, child anchor=west, anchor=west, edge={-Latex,draw=TreeLine}, l sep=9mm, s sep=1.2mm, font=\sffamily\fontsize{6.2}{7.2}\selectfont},
[{\textbf{3-catch\_ups/}\\\textcolor{DescText}{Remedial C5–C7 assessment material and its teacher solution companion.}}, rootnode, text width=6.0cm
[{\code{G10\_T3\_C5C6C7\_catch\_up.tex}\\\textcolor{DescText}{Student catch-up sheet on standard-normal probabilities, contextual standardization, and inverse-normal cutoffs/intervals.}}, filenode, text width=18.2cm]
[{\code{G10\_T3\_C5C6C7\_catch\_upsolutions.tex}\\\textcolor{DescText}{Teacher solutions for the catch-up sheet, with full calculations, distribution graphs, and a minimal-answer recap.}}, filenode, text width=18.2cm]
]
\end{forest}
\end{center}
\end{landscape}
\begin{landscape}
\subsection{\code{4-extra\_material/}}
\begin{center}
\begin{forest}
for tree={grow=east, parent anchor=east, child anchor=west, anchor=west, edge={-Latex,draw=TreeLine}, l sep=9mm, s sep=1.2mm, font=\sffamily\fontsize{6.2}{7.2}\selectfont},
[{\textbf{4-extra\_material/}\\\textcolor{DescText}{Standalone student reference material supplementing the normal-distribution work.}}, rootnode, text width=6.0cm
[{\code{normal-table.tex}\\\textcolor{DescText}{Printable standard-normal reference sheet with table-reading directions, identities, a shaded curve, and \ensuremath{\Phi} values through z = 3.09.}}, filenode, text width=18.2cm]
]
\end{forest}
\end{center}
\end{landscape}

\section{Resource roles}

\begin{itemize}
\item \textbf{Student-facing pages/materials:} the practice activities, context article, assessment sheets, catch-up sheet, project brief, datasets, and normal-table reference are intended for learners. In this repository they are printable LaTeX sources rather than web pages.
\item \textbf{Teacher materials:} files containing \code{solution}/\code{solutions} provide answer keys or worked derivations. \code{G10\_T3\_summary.tex} is explicitly a pedagogical review rather than student exercises.
\item \textbf{Lesson and assessment resources:} \code{1-practice\_activities/} develops criteria C0–C7 and also contains consolidated C1–C10 material; \code{2-learning\_evidences/} holds formal or practice assessment variants; \code{3-catch\_ups/} provides remediation.
\item \textbf{Stylesheets and JavaScript:} none are stored here. The C8–C10 brief requires students to create HTML, CSS, and JavaScript, but those deliverables are not part of this directory.
\item \textbf{Images and media:} no image file is stored here. Several assessment sources request a shared \code{logo.png} outside this subtree or through the shared LaTeX search path; graphs are generated in LaTeX with TikZ/PGFPlots or shared graph macros.
\item \textbf{Data and configuration:} the three CSV files are the only local datasets. There are no local configuration files.
\item \textbf{Downloadable documents:} no compiled PDF or office-document download is checked in; the \code{.tex} files are sources from which printable documents can be built.
\item \textbf{Generated/build files:} none are present or inventoried.
\end{itemize}

\section{Important relationships}

\begin{itemize}
\item Nearly every LaTeX document imports a shared preamble (\code{preamble\_exams.tex} or \code{preamble.tex}) using relative search paths. Those shared files live outside \code{term\_3} and supply document classes, packages, layout commands, and assessment helpers; the two \code{G10\_T3\_MATERIAL*.tex} files directly input \code{../../../preamble/preamble\_exams.tex}.
\item C4–C7 assessments and solution sets also import shared graph libraries such as \code{graphs/probability\_distribution\_graphs.tex}, \code{graphs/normal\_distribution\_graphs.tex}, and \code{graphs/uniform\_linear\_probability\_distributions.tex}. These external macro files generate reusable uniform, triangular, and normal plots; they are not assets contained in this tree.
\item Assessment headers in the midterm, C1–C4 exams, and catch-up sheet reference \code{logo.png} through the preamble/graphics search paths. The C8–C10 activity instead checks \code{../../../preamble/logo.png} explicitly and displays a placeholder if it is unavailable. No local copy of the logo exists.
\item \code{G10\_T3\_C1C7\_midterm\_solution.tex} corresponds to \code{G10\_T3\_C1C7\_midterm.tex}; their \code{V2} counterparts form the alternate pair. V2 changes actual model bounds, means/spreads, wording, and layout—not merely the filename.
\item \code{G10\_T3\_L1\_C4\_exam\_solution.tex} and \code{G10\_T3\_L1\_C4\_exam\_solution\_long.tex} both answer the original C4 exam: the former is a compact grouped key, while the latter gives full per-question explanations. \code{G10\_T3\_L1\_C4\_examV2\_solution.tex} is specifically paired with the numerically different V2 practice exam.
\item \code{G10\_T3\_C5C6C7\_catch\_upsolutions.tex} is the worked companion to \code{G10\_T3\_C5C6C7\_catch\_up.tex}; both cover the same three criteria, while the solution source additionally loads shared graph macros.
\item The C8–C10 activity instructs students to build a separate interactive webpage. Its C9 data-modeling mode is expected to accept all three adjacent CSV datasets, estimate a suitable uniform/triangular/normal model, and compare a fitted PDF with an empirical histogram. No resulting HTML, CSS, JavaScript, report, or deployed-site configuration is included here.
\item \code{normal-table.tex} supports the standard-normal work in C5–C7 as a standalone printable reference. Some practice and solution documents embed their own shorter \ensuremath{\Phi} tables, so they do not directly \code{input} this file.
\item \code{G10\_T3\_MATERIAL\_V2.tex} is a substantive pedagogical revision of \code{G10\_T3\_MATERIAL.tex}: it preserves the C1–C10 collection structure while changing early continuous-density solutions to geometry-based methods and replacing at least one nonlinear model with a piecewise-constant one.
\end{itemize}

\section{Verification notes}

\begin{itemize}
\item Every item could be identified confidently from its contents; \textbf{no entry requires manual verification}.
\item No transient or generated item was present in the tree. In particular, there were no \code{.DS\_Store}, editor-cache, compiler-auxiliary, temporary, dependency, compiled PDF, or LaTeX build-output files to exclude.
\end{itemize}
\end{document}
